% ===================================================================
%  HOTO Na 9 — Dutse da garin ruwan magani. --- Kurmi da farauta.
%
%  Traduction, faite en 2026, en haoussa d'aujourd'hui, sur l'ido de
%  texto/io. Regles de la colonne : voir 10-hoto-01.tex. Deux scenes,
%  deux numerotations. Un bloc marque « ¶2 » par ordo.py,
%  t09-c1-05-1 : une seule cle, deux alineas.
%
%  « dutse » EST LA PIERRE ET C'EST AUSSI LA MONTAGNE, ET TOUTE LA
%  PREMIERE SCENE EST BATIE AVEC CE SEUL MOT. Le Sahel n'a pas de
%  Pyrenees ; il a des inselbergs, des dalles et des cailloux, et un
%  mot qui les couvre tous. La colonne en tire la chaine entiere :
%  duwatsu les montagnes (1), ƙwanƙolin duwatsu les cimes (3),
%  dutsen wuta le volcan (6) — la montagne de feu —, dutsen ƙanƙara
%  le glacier (5) — la montagne de glace —, ma’adinan dutse la
%  carriere (30), dutsen yashi le gres. C'est le meme procede que la
%  neige du tableau 4 : quand la chose manque, on la compose avec ce
%  qu'on a, et l'on ne va rien chercher ailleurs. SEPTIEME PAIRE
%  D'HOMONYMES de cette colonne, apres kaka, kwano, allo, allura,
%  keke et ƙugiya — et la plus productive de toutes.
%
%  ET « keke » PARAIT POUR LA CINQUIEME FOIS : le funiculaire (13)
%  est « keken igiya », la machine a cable. Bicyclette, voiture
%  attelee, rouet, roue de moulin, funiculaire. Rien de commun entre
%  ces cinq objets, sinon un axe.
%
%  L'ARAIGNEE (1) EST « gizo-gizo » ET SA TOILE (2) EST « yanar
%  gizo », ET C'EST LE MEME GIZO QU'AU TABLEAU 8. L'arc-en-ciel y
%  etait « bakan gizo », l'arc de l'araignee : le fripon des contes
%  haoussa, celui qui ruse et gagne. Le voici en personne, entre deux
%  branches, a tendre son fil. MAIS « yanar gizo » EST AUSSI, EN
%  2026, LE NOM HAOUSSA DE L'INTERNET — les deux memes mots, sans une
%  lettre de changee, pour la toile tendue entre deux branches et
%  pour celle qui est tendue entre deux continents. La consigne est
%  d'ecrire la langue de 2026 ; ici elle n'a rien demande.
%
%  LE CYGNE (45) N'A PAS DE NOM, ET LA COLONNE JAVANAISE AVAIT EU LE
%  PROBLEME INVERSE. Il n'y a pas de cygne en Afrique ; on ecrit
%  « farin agwagwa mai doguwar wuya », le canard blanc a long cou. Le
%  javanais, lui, avait deux mots trop proches et s'etait trompe :
%  « banyak » l'oie pour « angsa » le cygne, corrige apres coup. Deux
%  colonnes, deux difficultes exactement contraires sur le meme
%  oiseau — trop de mots la, pas un seul ici.
%
%  SEPT ARBRES, SIX DECRITS ET UN EMPRUNTE. Le sapin (31) est
%  « itatuwa masu allura », les arbres a aiguilles ; le hetre (23),
%  celui a l'ecorce lisse ; le bouleau (26), celui a l'ecorce
%  blanche ; le pin, celui a la resine ; le cedre (49), celui qui
%  sent ; l'orme (47), l'arbre-parasol. Et le peuplier (48) recoit
%  MOT POUR MOT la description qu'il avait deja au tableau 3 —
%  « sirarun bishiyoyi » —, ce qui est le seul controle de coherence
%  qu'une colonne puisse s'appliquer a elle-meme. Seul le chene (29)
%  est emprunte, « itacen oak », parce que c'est le seul des sept
%  dont le BOIS se vend au Nigeria sous ce nom-la.
%
%  ET « fadama » REND « marsho » (44) SANS UN MOT DE TROP. La fadama
%  est la plaine inondable des vallees haoussa, celle que la crue
%  couvre et fertilise, et c'est la terre la plus chere du pays.
%  L'ido dit qu'une inondation recente avait change le chemin en vrai
%  marecage : le haoussa a le mot exact, et il y met en prime tout ce
%  que le mot vaut.
%
%  UNE FAUTE D'ORDRE, LA QUATRIEME, ET C'EST LA DEUXIEME FOIS QU'UNE
%  AUTRE COLONNE S'ETAIT TROMPEE AU MEME BLOC. Au bloc t09-c2-09-1 le
%  cerf (50) avait ete pose avant le taillis (54) d'ou il debouche,
%  alors que l'ido dit « me vidis debushar ek kopso belega cervulo » :
%  la source d'abord, l'animal ensuite. La colonne PERSANE avait fait
%  exactement cette faute, a ce bloc-ci, et la javanaise avait fait la
%  meme au bloc t03-c2-07-1 ou le haoussa a glisse aussi. TROIS
%  LANGUES SANS AUCUN RAPPORT — persan, javanais, haoussa — qui se
%  trompent aux deux memes endroits du livret : ce ne sont pas les
%  langues qui fautent, ce sont deux phrases de l'ido qui posent le
%  lieu avant l'objet la ou toute langue a tete initiale le pose
%  apres. Le remede est le meme : fronter la source, « daga
%  sarƙaƙƙiya... yana fitowa ».
%
%  L'OURS (44) REVIENT POUR LA TROISIEME FOIS, ET C'EST LA PREMIERE
%  OU CE SOIT LA BETE. Au tableau 2 c'etait un jeu, au tableau 7 une
%  enseigne d'auberge qui grincait au vent du nord ; ici le montreur
%  le tient par une chaine. Un mot emprunte pour un animal qu'aucun
%  Haoussa n'a vu aura servi trois fois avant de designer l'animal.
% ===================================================================

%%K t09-apar-1 apar
\VUsaut{4.0mm}
\VUtitre{15.0pt}{60}{HOTO Na 9}
\VUsaut{3.0mm}

%%K t09-tit-1 sub
\VUcentre{12.0pt}{0}{\textit{Fage na Farko.}}
\VUsaut{3.0mm}

%%K t09-tit-2 sub
\VUpk{11.4pt}{40}{Dutse. --- Garin ruwan magani.}
\VUsaut{3.0mm}

%%K t09-c1-01-1 p
1. --- Kwana takwas ke nan ina Pratz. Ba zan iya bayyana dukan
sha’awar da na ji lokacin da nake gaban \VUgras{duwatsu}
\textsuperscript{(1)} masu ban mamaki na Pirenei ba, waɗanda
\VUgras{kansu} \textsuperscript{(2)} yake bayyana a kan shuɗiyar sama kamar
babban ado.

%%K t09-c1-02-1 p
2. --- Ban zata cewa \VUgras{ƙwanƙolin duwatsu} \textsuperscript{(3)} na
Pirenei suna kai wa \VUgras{tsawo} mai girma haka ba, na kuma yi mamaki
ƙwarai gaban kyawawan \VUgras{duwatsun ƙanƙara} \textsuperscript{(5)} da
\VUgras{ƙwanƙolin dutse} \textsuperscript{(4)} masu dusar ƙanƙara waɗanda
\VUgras{zaftarewar ƙanƙara} mai ban tsoro take gudana a kansu.

%%K t09-tit-3 sub
\VUsaut{3.0mm}
\VUpk{10.2pt}{40}{Yanayin dutse.}
\VUsaut{3.0mm}

%%K t09-c1-03-1 p
3. --- Tun washegarin isowata, na je tsohon \VUgras{dutsen wuta}
\textsuperscript{(6)} da ya mutu, wanda yake kusa da Pratz. A gaɓoɓinsa
busasshe kuma marasa komai na \VUgras{bakin dutsen wuta,} har yanzu
ana ganin alamun wucewar \VUgras{narkakken dutse} sarai, wanda ya
gudana, ƙarnuka da yawa da suka wuce, a kan \VUgras{gefen dutsen}
\textsuperscript{(7)}. Na yi nasarar samun, a kan ɗaya daga cikin
\VUgras{gangarorin,} waɗansu guntayen wannan narkakken dutse.

%%K t09-c1-04-1 p
4. --- Pratz yana kan iyaka, \VUgras{kagara} \textsuperscript{(8)} kuwa
wadda take kare \VUgras{mashigin} \textsuperscript{(27)} da yake raba
Faransa da Sifaniya, tana tunatar da tsofaffin kagarorin gaɓar Rhine
gaba ɗaya, waɗanda suke kama da masu leƙon \VUgras{ganima} daga saman
dutsensu.

%%K t09-c1-05-1 p
5. --- Babbar \VUgras{mikiya} \textsuperscript{(9)} wadda take da sheƙarta
ba nesa da \VUgras{kagara} \textsuperscript{(8)} ba, sau da yawa takan ja
hankalina.

Wannan \VUgras{tsuntsu mai farauta,} wanda \VUgras{bakinsa} yake da
ƙarfi, sau da yawa yakan kawo wa ƴaƴansa ɗan tunkiya ko ɗan akuya
wanda yake riƙe da \VUgras{faratansa.} Girman \VUgras{fikafikansa} ya
kai ta yadda zai iya shawagi da tsawon lokaci tare da kayansa masu
nauyi.

%%K t09-tit-4 sub
\VUsaut{3.0mm}
\VUpk{10.2pt}{40}{Matafiyi.}
\VUsaut{3.0mm}

%%K t09-c1-06-1 p
6. --- A can yawon da ya fi daɗi shi ne tafiya zuwa faɗuwar ruwa
\textsuperscript{(10)} ta « Kau ». \VUgras{Faɗuwar ruwan} tana faɗuwa
tana juyawa sa’an nan ta ɓace ta zama kyakkyawan \VUgras{ƙaramin kogi}
a cikin \VUgras{kurmin itatuwa masu allura} \textsuperscript{(11)} da yake
yamma da garin. Mun hau ta wannan kurmi har zuwa \VUgras{wurin lura da
yanayi} \textsuperscript{(12)}, kuma daga saman \VUgras{wurin kallo,} mun
hangi \VUgras{duk yankin} na wurin. \VUgras{Yanayin ƙasa} yana da ban
mamaki, \VUgras{halitta} kuwa ta yi kama da wadda take ba wa ƙasar
dukan dukiyar da take da ita.

%%K t09-c1-07-1 p
7. --- Domin mu sauka, mun yi amfani da \VUgras{keken igiya}
\textsuperscript{(13)}, mun kuma tsaya a ƙaramar \VUgras{tasha} kusa da
\VUgras{kufan} \textsuperscript{(14)} tsohuwar \VUgras{kagara} wadda
sarakunan Pratz suka zauna a ciki a dā.

%%K t09-c1-08-1 p
8. --- Kagara mai tausayi ! Me take tunani kuwa tana ganin a ƙasa
\VUgras{garin ruwan magani} \textsuperscript{(15)} mai kwarjini wanda yanzu
ya zama \VUgras{tashar ruwan zafi} ta zamani ?

%%K t09-c1-08-2 p
\VUgras{Yanayi} yana da daɗi ƙwarai, \VUgras{ruwan ma’adinai} kuwa
yana da amfani ƙwarai, har \VUgras{maganin} da ake yi a
\VUgras{asibitin hutu} \textsuperscript{(17)} jin daɗi ne na gaskiya.

%%K t09-tit-5 sub
\VUsaut{3.0mm}
\VUpk{10.2pt}{40}{Garin ruwan zafi mai daɗi.}
\VUsaut{3.0mm}

%%K t09-c1-09-1 p
9. --- Ban da haka, an yi kome domin zaman ya yi daɗi. An gina babbar
\VUgras{gadar jirgi} \textsuperscript{(16)} domin a sami hanyar
jirgin ƙasa ; an shirya \VUgras{gonar shaƙatawa} \textsuperscript{(18)} ta
\VUgras{gidan wanka na magani,} inda ake \VUgras{kiɗa,} ta hanya mai
kwarjini. An gina « \VUgras{gidajen hutu} » \textsuperscript{(19)} masu
kyau da yawa kewaye da gidan caca \textsuperscript{(21)}, \VUgras{babbar
hanya} \textsuperscript{(22)} kuwa da ake kula da ita sosai tana sauƙaƙa
sadarwa ƙwarai.

%%K t09-c1-10-1 p
10. --- \VUgras{Gangaren ƙasa} \textsuperscript{(23)} kawai yake raba
\VUgras{hanya} \textsuperscript{(20)} da \VUgras{kwari} \textsuperscript{(25)}
na gaskiya, wanda a ƙarshensa akwai kyakkyawan ƙaramin \VUgras{tafki}
\textsuperscript{(26)} wanda yake ciyar da \VUgras{ƙaramin rafi}
\textsuperscript{(24)} mai tsabta.

%%K t09-c1-11-1 p
11. --- A ɗaya gefen kwarin, ana aikin \VUgras{ma’adinan ƙarfe}
\textsuperscript{(28)}. Sau da yawa na je na ga \VUgras{ma’aikatan
ma’adinai} suna sauka cikin \VUgras{rami,} har ma na shiga
\VUgras{hanyar ƙasa} inda suke ciyar da yininsu, suna fito da
\VUgras{ma’adini} wanda yake ɗauke da \VUgras{ƙarfe.}

%%K t09-c1-12-1 p
12. --- An gina babban \VUgras{wurin narkar da ƙarfe} kusa da
ma’adinan domin a yi amfani da dukiyar da ake fito da ita nan take.
\VUgras{Babbar tanda} \textsuperscript{(29)}, wadda ake ɗumama ta da
\VUgras{gawayin dutse} wanda yake da yawa a nan, tana ba da damar samun
\VUgras{narkakken ƙarfe} da sauri kuma da arha.

%%K t09-c1-13-1 p
13. --- Ba nesa da ma’adinan ba ana haƙa \VUgras{ma’adinan dutse}
\textsuperscript{(30)}. Ba \VUgras{dutsen granit,} ba \VUgras{dutsen
forfir,} ba ma \VUgras{dutsen yashi} tsohon \VUgras{mai haƙar dutse}
yake yankewa da \VUgras{gatarin dutsensa,} sai \VUgras{sassaƙaƙƙun
duwatsu} kaɗai domin ginin gidaje.

%%K t09-c1-14-1 p
14. --- Ɗaya daga cikin kyawawan adon wannan ƙasa shi ne
\VUgras{itatuwa masu allura} \textsuperscript{(31)}. Ko’ina a ƙarshen
\VUgras{kwazazzabo} \textsuperscript{(32)}, a gaɓar \VUgras{gudun ruwa}
\textsuperscript{(33)}, ko \VUgras{rami mai zurfi} \textsuperscript{(34)} ko
gaci, sirarun allurarsu suna ba da inuwar kore.

%%K t09-tit-6 sub
\VUsaut{3.0mm}
\VUpk{10.2pt}{40}{Tafiya da ƙafa.}
\VUsaut{3.0mm}

%%K t09-c1-15-1 p
15. --- Ina amfani da hutuna domin \VUgras{doguwar} tafiya.
\VUgras{Hanyoyi} \textsuperscript{(35)} da suke karkata a kan dutsen suna
ba da damar tafiya, ba tare da lura da nisa ba, \VUgras{kilomita} da
yawa kuma sau da yawa \VUgras{kilomita huɗu-huɗu} da yawa.

%%K t09-c1-15-2 p
\VUgras{Iyakoki} \textsuperscript{(36)} da \VUgras{sandunan nuni}
\textsuperscript{(37)} suna nuna hanyoyin, inda sau da yawa mutum yakan
haɗu da matashin \VUgras{ɗan dutse} \textsuperscript{(38)} da \VUgras{jaka
mai kai biyu} \textsuperscript{(40)} a gefensa, yana ɗauke da
\VUgras{ɓeran dutse} \textsuperscript{(39)}, ko wani \VUgras{ɗan yawo}
\textsuperscript{(41)} wanda \VUgras{gajeriyar rigarsa}
\textsuperscript{(42)} take yin bambanci mai ban mamaki da tsohuwar
\VUgras{hular kansa} \textsuperscript{(43)} da ta yage. Yana jan
\VUgras{beyar} \textsuperscript{(44)} da aka hora da sarƙa.

%%K t09-tit-7 sub
\VUpk{10.2pt}{40}{Hawan duwatsu.}
\VUsaut{3.0mm}

%%K t09-c1-16-1 p
16. --- Kusan kowace rana ina tafiya tare da \VUgras{jagora}
\textsuperscript{(48)} domin in yi \VUgras{hawan dutse,} ina kuma
alfahari ƙwarai idan, da \VUgras{jaka} \textsuperscript{(47)} a baya da
« \VUgras{sandar hawa} » a hannu, kamar \VUgras{matafiyi}
\textsuperscript{(46)} na gaskiya, na kai wa \VUgras{duwatsu}
\textsuperscript{(45)} hari.

%%K t09-c1-17-1 p
17. --- Daga saman dutse, mazaunan fili suna kama da waɗanda ba su fi
tururuwa girma ba ; \VUgras{garken tumaki} \textsuperscript{(50)} yana
kiwo a \VUgras{makiyaya} \textsuperscript{(49)} ya yi kama da abin wasan
yara. \VUgras{Itacen ɗanko} \textsuperscript{(51)} ko kuma \VUgras{ƙaramin
gidan katako} \textsuperscript{(52)} yana kama da ɗigo kaɗai a kan
kore.

%%K t09-c1-18-1 p
18. --- A lokacin waɗannan tafiye-tafiye, sau da yawa mukan haɗu a kan
\VUgras{ƙaramar hanya} \textsuperscript{(53)} da \VUgras{mafarautan
akuyar dutse} \textsuperscript{(54)} yana ɗauke da dabbar da ya kashe yanzu
a kafaɗa.

%%K t09-c1-19-1 p
19. --- Na ajiye a littafin ganyayena reshen \VUgras{ganyen zare}
\textsuperscript{(55)} mai ban sha’awa ƙwarai da kyakkyawan ƙaramin
reshe mai launin hoda na \VUgras{furen tudu} \textsuperscript{(57)}
waɗanda na tsinka a bakin ƙaramin \VUgras{kogo} \textsuperscript{(56)} a
tafiyata ta ƙarshe.

%%K t09-tit-8 sub
\VUsaut{3.0mm}
\VUcentre{12.0pt}{0}{\textit{Fage na Biyu.}}
\VUsaut{3.0mm}

%%K t09-tit-9 sub
\VUpk{11.4pt}{40}{Kurmi. --- Farauta.}
\VUsaut{3.0mm}

%%K t09-c2-01-1 p
1. --- Jiya na yi rana mai daɗi ƙwarai a cikin kurmi. Himmar da take
tsakanin dabbobi da \VUgras{ƙwari} \textsuperscript{(5)} da suke zaune a
ciki ta ba ni sha’awa ƙwarai.

%%K t09-c2-01-2 p
\VUgras{Gizo-gizo} \textsuperscript{(1)} wanda yake saƙa \VUgras{yanarsa}
\textsuperscript{(2)} tsakanin rassa biyu domin ya kama \VUgras{ƙuda}
\textsuperscript{(3)} mai wucewa, \VUgras{katantanwa}
\textsuperscript{(4)} tana ɗauke da ƙwaryarta da wahala, ƙwaro mai ƙaho
yana bincike a kan abin da ya kama, tururuwa mai himma tana ƙwazo kusa
da \VUgras{gidan tururuwa} \textsuperscript{(6)}, dukan waɗannan ƙanana
sun tafi, sun zo, sun yi aiki da himma ta musamman.

%%K t09-c2-02-1 p
2. --- \VUgras{Maciji} \textsuperscript{(7)} wanda da farko na zata
\VUgras{kububuwa,} amma wanda ba wani abu ba ne sai
\VUgras{maciji marar dafi} kaɗai, ya ɓoye a ƙarƙashin kututturen
bishiya, lokaci zuwa lokaci kuma yakan fito da ƙaramin kansa mai
sirara wanda kullum ya yi kama da mai shirin cizo. A gabana
\VUgras{babbar katantanwa} \textsuperscript{(8)} tana rarrafe da nauyi
bayan ta cinye ɗaya daga cikin \VUgras{ƴaƴan sitiroberi}
\textsuperscript{(11)} masu daɗi da suke girma da yawa a can.

%%K t09-c2-03-1 p
3. --- An shimfiɗa ƙasar da \VUgras{furannin kurmi} ;
\VUgras{furannin farko} \textsuperscript{(9)}, da \VUgras{furen shuɗin
daji} \textsuperscript{(10)}, da waɗansu tsire-tsire da yawa da ban sani
ba, suna yin fure tsakanin \VUgras{kunkuman ciyawa}
\textsuperscript{(12)}.

%%K t09-c2-04-1 p
4. --- A wannan yanki, \VUgras{naman kaza na ƙasa} kusan yana da yawa
kamar \VUgras{naman kaza} \textsuperscript{(14)}, kuma \VUgras{ɗan alade}
\textsuperscript{(13)} na wani mai gadin daji, yana bincike da haɗama ko
zai yi nasarar samun waɗansu daga cikinsu.

%%K t09-tit-10 sub
\VUsaut{3.0mm}
\VUpk{10.2pt}{40}{Farautar sata.}
\VUsaut{3.0mm}

%%K t09-c2-05-1 p
5. --- Na halarci \VUgras{ƙarshen} wani fage da ya ba ni daɗi ƙwarai.
Da taimakon sunsunar \VUgras{ƙaramin karen farautarsa}
\textsuperscript{(15)}, \VUgras{mai gadin daji} \textsuperscript{(16)} ya
kama \VUgras{ɗan farautar sata} \textsuperscript{(22)} bisa ba zato ba
tsammani, a daidai lokacin da wannan yake shirin ɗaukar \VUgras{naman
daji} \textsuperscript{(17)} da ya kashe. Domin ya fi son ya bar abin da
ya kama da a kama shi, ɗan farautar satan ya gudu, \VUgras{zomon daji}
\textsuperscript{(18)}, da \VUgras{barewa} \textsuperscript{(19)}, da
\VUgras{ƙaramar barewa} \textsuperscript{(20)} da \VUgras{aladen daji}
\textsuperscript{(21)} kuwa suka kwanta a kan ciyawa suna faranta wa mai
gadin daji rai.

%%K t09-c2-06-1 p
6. --- Bambancin itatuwan da kurmi yake ɗauke da su yana da ban
mamaki. \VUgras{Itatuwa masu santsin ɓawo} \textsuperscript{(23)} da
\VUgras{itatuwa masu farin ɓawo} \textsuperscript{(26)}, da
\VUgras{itatuwa masu ɗanko,} waɗanda suke ba da \VUgras{ɗanko}
\textsuperscript{(27)}, da \VUgras{itatuwan oak} \textsuperscript{(29)} da
kuma waɗansu da yawa suna haɗa rassansu.

%%K t09-c2-06-2 p
\VUgras{Ciyawar hawa} \textsuperscript{(24)}, wadda take manne da
kututturen dogayen itatuwa, tana ba wa \VUgras{kurmi}
\textsuperscript{(25)} launin kore mai ƙyalli, wanda yake haɗuwa da daɗi
da launin kore mai haske na \VUgras{gansakuka} \textsuperscript{(31)}
inda \VUgras{tsuntsun ƙwanƙwasa itace} \textsuperscript{(30)} yake neman
ƙwari.

%%K t09-tit-11 sub
\VUsaut{3.0mm}
\VUpk{10.2pt}{40}{Yanayin kurmi.}
\VUsaut{3.0mm}

%%K t09-c2-07-1 p
7. --- Tsofaffin itatuwan oak sun ba ni sha’awa. Manyan
\VUgras{saiwoyinsu} \textsuperscript{(32)} da suka murɗe, rabin fitowa
daga ƙasa, suna ba su kyakkyawan kamannin ƙarfi da ƙwari, ina kuma
tambayar kaina yadda \VUgras{mai gonar} \textsuperscript{(35)} zai iya
yarda a sassare su a kai su ga \VUgras{zarton} \textsuperscript{(34)}
\VUgras{mai sara} \textsuperscript{(33)}. \VUgras{Kututture}
\textsuperscript{(36)} masu tausayi, da aka tuɓe \VUgras{ɓawonsu}
\textsuperscript{(37)} ko da \VUgras{gatarin} \textsuperscript{(38)}
\VUgras{ɗan ƙwadago} \textsuperscript{(39)} ya tsage su, suna ɓata mini
rai.

%%K t09-c2-08-1 p
8. --- A kusa, tsoho \VUgras{mai yin gawayi} \textsuperscript{(41)} ya
shirya da \VUgras{shebur ɗinsa tudun} \textsuperscript{(42)} gawayi,
tsohuwa kuwa ta ɗauko \VUgras{ɗamin} \textsuperscript{(40)}
\VUgras{busasshen itace} na rassan itatuwan da aka sassare.

%%K t09-tit-12 sub
\VUsaut{3.0mm}
\VUpk{10.2pt}{40}{Farauta.}
\VUsaut{3.0mm}

%%K t09-c2-09-1 p
9. --- Na so in je in huta na ɗan lokaci a \VUgras{gidan mai gadin
daji} \textsuperscript{(46)} amma, lokacin da na iso kusa da \VUgras{ƙaramin
tafki} \textsuperscript{(43)}, inda kyakkyawar \VUgras{farin agwagwa mai
doguwar wuya} \textsuperscript{(45)} take iyo, sai na lura cewa
\VUgras{ambaliya} \textsuperscript{(44)} ta baya-bayan nan ta canja hanyar
ta zama \VUgras{fadama} ta gaskiya. Sai na juya, kuma, ina wucewa kusa
da \VUgras{itacen ƙamshi} \textsuperscript{(49)}, da \VUgras{sirarun
bishiyoyi} \textsuperscript{(48)} da \VUgras{itacen laima}
\textsuperscript{(47)} da suke a bakin kurmin, sai na ga daga
\VUgras{sarƙaƙƙiya} \textsuperscript{(54)} \VUgras{namijin barewa}
\textsuperscript{(50)} mai kyau yana fitowa, \VUgras{karnukan farauta}
\textsuperscript{(51)} masu taurin kai suna korarsa, waɗanda
\VUgras{ƙahon} \textsuperscript{(53)} \VUgras{mafarauta a kan dawakai}
\textsuperscript{(52)} ya ƙarfafa su. Dabbar mai tausayi ta gaji gaba
ɗaya. Ta zo ta faɗi tana mutuwa ƴan taku kusa da ni.

%%K t09-c2-10-1 p
10. --- Ɗan mai gadin daji, wanda hayaniyar \VUgras{farautar gudu} ta
jawo shi, ya fito daga gidan mu kuma mu biyu muka koma cikin kurmin.
Ya ga \VUgras{tsuntsu mai baƙi da fari} \textsuperscript{(56)} a kan
reshen tsohon itacen oak. Ya zaci cewa wataƙila an ɓoye sheƙarsa cikin
\VUgras{ganyaye} \textsuperscript{(58)}, ya kuma so ya kama sheƙar.

%%K t09-c2-10-2 p
Mai saurin motsi kamar \VUgras{kurege} \textsuperscript{(61)}, ya hau
daga \VUgras{reshe} \textsuperscript{(55)} zuwa reshe, amma bai sami komai
ba ya kuma yi nasarar tashi da tsohuwar \VUgras{hankaka}
\textsuperscript{(60)} kaɗai wadda take zaune a kan \VUgras{ƙaton reshe}
\textsuperscript{(57)}.

\VUsaut{4.0mm}
\VUfilet{22mm}
